% Episode IV — All Eyes on You (Cat vs Dog)
% Calibré sur la structure et le style de l'Episode III — Birth of a Neuron
%
% Si TeX Live complet : [most]{tcolorbox}, microtype, tabularx+ragged2e

\documentclass[11pt]{article}
\usepackage[utf8]{inputenc}
\usepackage[T1]{fontenc}
\usepackage{tcolorbox}
\usepackage{amsmath, amssymb}
\usepackage{tikz}
\usepackage{graphicx}
\usepackage{listings}
\usepackage{enumitem}
\usepackage{hyperref}
\usepackage{lmodern}
\usepackage{xcolor}
\usepackage{caption}
\usepackage{float}
\usepackage{fancyvrb}
\usepackage[margin=1in]{geometry}
\usepackage{xcolor}
\usetikzlibrary{arrows.meta, positioning}

\hypersetup{
  colorlinks=true,
  linkcolor=blue,
  urlcolor=blue,
  citecolor=blue
}

% --- Code style (minimal, aligné Ep III) ---
\definecolor{backcolour}{rgb}{0.95,0.95,0.92}
\lstdefinestyle{mystyle}{
    backgroundcolor=\color{backcolour},
    basicstyle=\ttfamily\footnotesize,
    breaklines=true
}
\lstset{style=mystyle}

% --- Chemins images ---
\graphicspath{{../../assets/photos_git/}}

% ---------------------------------------------------------------
%                      TITLE PAGE
% ---------------------------------------------------------------
\begin{document}
\begin{titlepage}
    \centering
    \vspace*{2.2cm}

    {\Large \textbf{DEEP LEARNING FROM SCRATCH}}\\[0.4em]
    A Guide for Self-Learners\\[2.5em]

    {\Huge \bfseries Episode IV}\\[0.4em]
    {\LARGE \textit{All Eyes on You}}\\[2em]

    \vspace{2cm}

   \begin{tcolorbox}[colback=black!3!white, colframe=black!20!white, width=0.8\textwidth, sharp corners, boxrule=0.4pt]
    \centering
    \itshape

    {\Large Speed isn't the enemy, confusion is.}\\[1em]

    \textbf{Why this mini-guide?}\\[0.3em]
    We open its eyes.\\
    Until now, the neuron knew only numbers. Now we give it sight.\\
    For the first time, it will see — and yes, it will stumble.\\
    But here, we awaken that power.
    \end{tcolorbox}
    \vspace{1.5cm}
    \begin{center}
    \itshape
        {\Large \textit{"The most difficult thing is the decision to act; the rest is merely tenacity."}}\\
        \vspace{0.2em}
        {\large -- Amelia Earhart}
    \end{center}


    \vfill

    {\large \textbf{Pierre Chambet}}\\[0.3em]
    {\small February 2026}\\[1em]
    \rule{0.8\textwidth}{0.4pt}\\[0.4em]
    {\small Part of the \textit{Deep Learning From Scratch} series — an independent reconstruction of AI.}

\end{titlepage}


\begin{tcolorbox}[colback=white, colframe=black!80, boxrule=0.8pt,
  width=\textwidth, before skip=10pt, after skip=10pt, halign=center,
  fontupper=\sffamily\large]

\href{https://colab.research.google.com/github/Pchambet/Deep-Learning-from-Scratch/blob/main/notebooks/04_training_loop_from_scratch.ipynb}{
  \textbf{\large Launch All Eyes on You here on Google Colab}
}

\vspace{0.5em}

{\color{gray!80!black}\small%
\textit{One click. Your turn to run it.}
}
\end{tcolorbox}

\section*{Here, You Test Our Neuron on Real Images}

\subsection*{From Plants to Pixels}

In Episode III, we built a neuron that classified plants using two features: leaf length and leaf width.  
Now we ask: \textbf{can the same neuron tackle real images?}

We will use a small dataset of photos — \textbf{cats} and \textbf{dogs} — and see if the neuron can learn to tell them apart.  
Same equations. Same learning loop. The only change: each input has \textbf{4096 numbers} instead of 2.

\begin{tcolorbox}[colback=gray!5, colframe=gray!60!black]
\textbf{You don't need new theory.}  
You already know how the neuron works.  
What changes is the \emph{shape} of the input — and how we prepare it.
\end{tcolorbox}

\subsection*{Theory in Short (Reminder)}

Nothing changes in the neuron itself.  
We still have:
\begin{enumerate}
    \item \textbf{Model} — linear transformation $Z = XW + b$, then sigmoid.
    \item \textbf{Loss} — log-loss measures how wrong the predictions are.
    \item \textbf{Gradients} — derivatives tell us how to adjust $W$ and $b$.
    \item \textbf{Update} — gradient descent: $W \leftarrow W - \alpha \, dW$, $b \leftarrow b - \alpha \, db$.
\end{enumerate}

\begin{quote}
The neuron does not care whether $X$ has 2 columns or 4096.  
It applies the same math. Our job is to turn a photo into a vector of numbers.
\end{quote}


\begin{tcolorbox}[colback=gray!3, colframe=black!40, title=Let's test our neuron on images, fonttitle=\bfseries, arc=2mm]
We will load cats and dogs, flatten each image into 4096 pixels, normalize them, and feed them to the neuron.  
Along the way, we will fix a numerical bug, tune hyperparameters, and discover overfitting.  
\textbf{No magic — just the same loop, scaled up.}
\end{tcolorbox}

\vspace{0.7em}

\begin{center}
\textit{First question every data scientist asks: \textbf{what is my dataset?}}

Because behind every dataset, there is a goal.  
Here: can a single neuron learn to distinguish cats from dogs?
\end{center}

\vspace{0.7em}

\subsection*{The Dataset — Cats and Dogs}

\vspace{0.5em}

We use two separate datasets:
\begin{itemize}
    \item One for \textbf{training} — the neuron learns from it.
    \item One for \textbf{testing} — we evaluate on data the neuron has never seen.
\end{itemize}

\vspace{0.5em}

\noindent Why two?

\begin{tcolorbox}[colback=white!98!gray, colframe=gray!30!black, arc=1mm]
If we trained and tested on the same data, the neuron might \emph{memorize} answers without learning to generalize.  
By testing on \textbf{unseen images}, we check if it truly \textbf{understands} the task.
\end{tcolorbox}

\vspace{0.5em}

\noindent It is like preparing for an exam: practicing helps, but the real test is solving \textbf{new problems}.

\vspace{1em}

\subsection*{Step 1: Load and Inspect}

\begin{tcolorbox}[colback=gray!3, colframe=black!40, title=Step 1: Load the Neuron and the Data, fonttitle=\bfseries, arc=2mm]
We reuse the code from Episode III and load the cat/dog dataset.
\end{tcolorbox}

\vspace{0.5em}

\noindent In Colab, you'll import the neuron from Episode III and load the data. One line:

\begin{lstlisting}[language=Python]
X_train, y_train, X_test, y_test = load_data()
\end{lstlisting}

\vspace{0.5em}

\noindent To inspect what we have — \texttt{print(X\_train.shape)} — and we see:

\begin{tcolorbox}[colback=gray!5, colframe=black!40, arc=1mm, title=Training Data Overview, fonttitle=\bfseries]
    \begin{itemize}
        \item \texttt{X\_train shape}: \texttt{(1000, 64, 64)} — 1000 images, each 64×64 pixels.
        \item \texttt{y\_train shape}: \texttt{(1000, 1)} — one label per image.
        \item Labels: \texttt{0} = cat, \texttt{1} = dog. Balanced: 500 cats, 500 dogs.
\end{itemize}
\end{tcolorbox}

\vspace{0.5em}

\noindent Same idea for the test set:

\begin{tcolorbox}[colback=gray!5, colframe=black!40, arc=1mm, title=Test Data Overview, fonttitle=\bfseries]
    \begin{itemize}
        \item \texttt{X\_test shape}: \texttt{(200, 64, 64)} — 200 unseen images.
        \item \texttt{y\_test shape}: \texttt{(200, 1)}.
        \item Balanced: 100 cats, 100 dogs.
    \end{itemize}
\end{tcolorbox}

\begin{figure}[H]
    \centering
    \includegraphics[width=0.6\textwidth]{dogs}
    \caption{Sample images from the dataset.}
\end{figure}

\begin{tcolorbox}[colback=white!97!gray, colframe=black!20, arc=1mm]
\textbf{Pro tip:} Always inspect the shapes before modeling.  
You want to know exactly what you are feeding to the neuron.
\end{tcolorbox}

\vspace{0.5em}

\noindent The full setup — imports, paths — is in Colab. Nothing to worry about.

\vspace{1em}

\section*{How Do We Feed a Photo to the Neuron?}

\vspace{0.5em}

\noindent Now: how do we actually give a photo to the neuron?

\noindent\textbf{The answer is simple.}  
We treat each pixel as a separate input.  
A 64×64 image has $64 \times 64 = \boldsymbol{4096}$ pixels.

\vspace{0.5em}

We scan the image row by row — left to right, top to bottom — and flatten it into a vector:
\[
\boldsymbol{X = (x_1, x_2, \ldots, x_{4096})}
\]

Each $x_i$ is the intensity of one pixel.

\vspace{1em}

\begin{tcolorbox}[colback=white!98!gray, colframe=gray!30!black, arc=1mm]
\textbf{A photo becomes a vector of numbers.}  
From there, the neuron applies the same linear model as before.
\end{tcolorbox}

\vspace{1em}

\begin{center}
    \fbox{\large \textbf{The neuron doesn't see images — it sees numbers.}}
\end{center}

\vspace{1em}

\section*{The Pipeline: Four Steps}

\vspace{0.5em}

How do we go from raw pixels to a prediction?

\begin{enumerate}
    \item \textbf{Normalize} — rescale each pixel from $[0, 255]$ to $[0, 1]$ (divide by 255 or by max).
    \item \textbf{Flatten} — turn each 64×64 image into a 4096-dimensional vector.
    \item \textbf{Train} — same loop as Episode III: model $\to$ loss $\to$ gradients $\to$ update.
    \item \textbf{Evaluate} — test on unseen data to check generalization.
\end{enumerate}

\begin{tcolorbox}[colback=white!98!gray, colframe=gray!30!black, arc=1mm]
\textbf{That is the pipeline.}  
From raw pixels to prediction — four steps.
\end{tcolorbox}

\vspace{1em}

\section*{Step 2: Normalize and Flatten}

\begin{tcolorbox}[colback=gray!3, colframe=black!40, title=Step 2: Prepare the Input, fonttitle=\bfseries, arc=2mm]
Before feeding images to the neuron, we flatten and normalize.  
Two operations, often done together.
\end{tcolorbox}

\vspace{0.5em}

\begin{itemize}
    \item \textbf{Flatten} — each 64×64 image becomes a vector of 4096 values.
    \item \textbf{Normalize} — scale pixels from $[0, 255]$ to $[0, 1]$.
\end{itemize}

Why divide by 255? Each pixel is stored on 8 bits: $2^8 = 256$ values, from 0 to 255.  
We divide by the max value — that way it works for any image format.

\begin{tcolorbox}[colback=white!97!gray, colframe=black!20, arc=1mm]
\textbf{Why normalize?}  
If one feature has much larger values than others, the model may over-weight it.  
Normalization ensures every feature contributes fairly — like balancing the mix in a band.
\end{tcolorbox}

\begin{center}
    \fbox{\textbf{We normalize to be fair.}}
\end{center}

\vspace{1em}

\noindent In code — two lines:

\begin{lstlisting}[language=Python]
X_train_reshape = X_train.reshape(X_train.shape[0], -1) / X_train.max()
X_test_reshape = X_test.reshape(X_test.shape[0], -1) / X_test.max()
\end{lstlisting}

\vspace{0.5em}

\begin{tcolorbox}[colback=gray!5, colframe=black!50, arc=2mm, title=At This Point..., fonttitle=\bfseries]
Each image is a normalized vector of 4096 values.  
The neuron can now learn from the data.
\end{tcolorbox}

\begin{tcolorbox}[colback=white!97!gray, colframe=gray!15!gray, arc=1mm]
\textbf{Premium tip:} Use \texttt{-1} in \texttt{reshape} to let NumPy infer the dimension.  
\texttt{X.shape[0], -1} means ``keep the number of images, flatten the rest.''
\end{tcolorbox}

\vspace{1em}

\noindent Data ready. Time to train.

\vspace{0.5em}

\section*{Step 3: Train — And Hit a Bug}

\begin{tcolorbox}[colback=gray!3, colframe=black!40, title=Step 3: Run the Training Loop, fonttitle=\bfseries, arc=2mm]
We feed the reshaped data to the neuron — same \texttt{artificial\_neuron} as in Episode III.
\end{tcolorbox}

\vspace{0.5em}

\begin{lstlisting}[language=Python]
W, b = artificial_neuron(X_train_reshape, y_train, learning_rate=0.1, num_iter=100)
\end{lstlisting}

\begin{figure}[H]
    \centering
    \includegraphics[width=0.6\textwidth]{error}
    \caption{The neuron's first attempt at seeing — and it stumbles.}
\end{figure}

\vspace{0.5em}

\noindent We get:
\begin{center}
    \emph{``RuntimeWarning: divide by zero encountered in log''}
\end{center}

\vspace{0.5em}

\subsection*{What is going wrong?}

The error comes from \texttt{log\_loss}.  
We compute $\log(A)$ and $\log(1-A)$.  
$A$ comes from the sigmoid: $A = 1/(1+e^{-Z})$.

\begin{tcolorbox}[colback=white!98!gray, colframe=gray!30!black, arc=1mm]
\textbf{The exponential grows fast.}  
For $z = 20$, we have $\exp(-z) \approx 0$, so $A \approx 1$.  
For $z = -20$, we have $A \approx 0$.  
And $\log(0)$ is undefined.
\end{tcolorbox}

\vspace{0.5em}

\noindent\textbf{Fix:} Add a small constant $\varepsilon$ to avoid $\log(0)$.  
Standard practice in deep learning.

\begin{lstlisting}[language=Python]
def log_loss(A, y):
    m = len(y)
    epsilon = 1e-15
    loss = -1/m * np.sum(y * np.log(A + epsilon) 
                    + (1 - y) * np.log(1 - A + epsilon))
    return loss
\end{lstlisting}

\begin{tcolorbox}[colback=white!97!gray, colframe=black!20, arc=1mm]
\textbf{Pro tip:} Numerical stability matters.  
Always guard against $\log(0)$, division by zero, and overflow.  
$\varepsilon = 10^{-15}$ is tiny — it does not change the math, but it prevents crashes.
\end{tcolorbox}

\vspace{0.5em}

\noindent Run the corrected loop in Colab and watch the curve stabilize.

\vspace{1em}

\section*{Step 4: Hyperparameters — The Levers We Can Turn}

\begin{tcolorbox}[colback=gray!3, colframe=black!40, title=Step 4: Tune the Learning Process, fonttitle=\bfseries, arc=2mm]
Once the model runs, we control how it learns via \textbf{hyperparameters}:  
learning rate, number of iterations, initialization, etc.
\end{tcolorbox}

\vspace{0.5em}

With the fixed \texttt{log\_loss}, training runs — but the curve may be shaky.

\begin{itemize}
    \item \textbf{Learning rate} — too high $\Rightarrow$ overshoot, instability.  
    Too low $\Rightarrow$ slow convergence.
    \item \textbf{Number of iterations (epochs)} — too few $\Rightarrow$ underfitting.  
    Too many $\Rightarrow$ risk of overfitting.
\end{itemize}

\begin{tcolorbox}[colback=white!98!gray, colframe=gray!30!black, arc=1mm]
\textbf{Hyperparameters = the teacher's schedule.}  
They decide how fast the neuron moves and how long it trains.
\end{tcolorbox}

\vspace{0.5em}

\noindent In Colab, we lower the learning rate and rerun. The curve smooths out.

\begin{figure}[H]
    \centering
    \includegraphics[width=0.6\textwidth]{better}
    \caption{With a gentler learning rate, the curve smooths out.}
\end{figure}

\vspace{0.5em}

We can also track both loss and accuracy — and compare training vs.\ test performance. In Colab, we run a loop that plots them side by side. What we see: the curves diverge.

\begin{figure}[H]
    \centering
    \includegraphics[width=1\textwidth]{definitely}
    \caption{Train climbs. Test lags. The gap tells the story.}
\end{figure}

\vspace{1em}

\section*{Step 5: Detecting Overfitting}

\begin{tcolorbox}[colback=gray!5, colframe=black!50, title=Step 5: Diagnose the Divergence, fonttitle=\bfseries, arc=2mm]
The plots reveal a clear pattern:  
training loss decreases, training accuracy climbs — but test loss stagnates or rises, test accuracy lags.
\end{tcolorbox}

\vspace{0.5em}

\noindent Look at the plot above — that gap is the signal. \textbf{Overfitting} means the model performs very well on training data but fails to generalize to new examples.

\begin{tcolorbox}[colback=white!98!gray, colframe=gray!30!black, arc=1mm]
Like a student who memorizes the exercises without understanding the method —  
they ace the practice set but struggle on new problems.
\end{tcolorbox}

\vspace{0.5em}

\begin{center}
    \fbox{\textbf{Warning: if train and test curves diverge, you are likely overfitting.}}
\end{center}

\vspace{1em}

\section*{How Can We Fix Overfitting?}

Common strategies:
\begin{enumerate}
    \item \textbf{More data} — expose the model to greater variety.
    \item \textbf{Simplify the input} — reduce dimensions, extract better features.
    \item \textbf{Regularization} — L1, L2, dropout — push weights toward simpler solutions.
\end{enumerate}

\vspace{0.5em}

But here's the twist: in our case, none of these will dramatically help. The problem isn't just overfitting. Why?

\begin{tcolorbox}[colback=white!98!gray, colframe=gray!30!black, arc=1mm]
\textbf{Because our model is too simple.}  
We use a \textbf{single neuron} — a linear classifier.  
It can only draw one straight line.
\end{tcolorbox}

\vspace{0.5em}

If the true boundary between cats and dogs is curved, complex, nonlinear —  
a single neuron is not enough.

\vspace{1em}

\section*{The Limit — and the Next Step}

\vspace{0.5em}

We have pushed our neuron as far as it can go.

It learned. It adjusted. It gave us decent results on the training set.  
But it hits a wall: one straight line cannot capture the complexity of images.

\begin{tcolorbox}[colback=gray!5, colframe=black!50, title=What do we do?]
\textbf{We scale up.}  
We add more neurons. We add more layers.  
We move from a single-cell brain to a full \textbf{Artificial Neural Network (ANN)}.
\end{tcolorbox}

\vspace{1em}

\begin{center}
    \fbox{\Large \textbf{Let's move from a neuron... to a brain.}}
\end{center}

\vspace{0.5em}

And yes — that means more equations, matrices, and weight updates.  
But you only need to understand it once.  
Then you will be free to design your own architectures.

\vspace{1em}

We gave it eyes. It saw. It stumbled. And now we know why — and what comes next.

\vspace{0.5em}

In the next episode, we build our first real neural network.

\end{document}
